% Chapter 1.
% PLACEHOLDER PROSE from the layout specimens. The facts are from memory and
% have NOT been checked against sources. Replace or verify before keeping.
\chapter{The Berry}

A grape is a berry. Botanists mean something exact by that word: a fleshy
fruit that grows from a single flower with a single ovary, its seeds buried
in the pulp. By this rule the tomato is a berry, and so are the banana and
the eggplant, while the strawberry and the raspberry are
not.\anchor{not-berries}\note[strawberry]{The strawberry is an accessory fruit. Most of its flesh grows from
the receptacle, the part of the flower that holds the ovaries, and the specks
on the outside are the true fruits, each one an
achene.\subnote{An achene is a small dry fruit with one seed that does not
split open when ripe. A sunflower seed in its striped shell is also an
achene, although botanists give the sunflower's version its own name,
cypsela.} The raspberry is an aggregate: a cluster of many small fruits, each
from a separate ovary of one flower. The strawberry comes back later in the
chapter, \xref{strawberry-returns}.} The grape qualifies without argument.

Nearly all the grapes we eat, dry, and press come from one species,
\emph{Vitis vinifera}, the wine-bearing
vine.\aside{\emph{Vitis} is the Latin for vine.}\note{\emph{Vinifera} is
Latin for wine-bearing, from \emph{vinum}, wine, and \emph{ferre}, to bear.
The English word \emph{grape} took a stranger road. It comes from Old French
\emph{grape}, a bunch of grapes, and probably before that from a Germanic
word for a hook.\subnote{Presumably the hook used to cut the bunches. The
same root is thought to lie behind \emph{grapple} and French
\emph{grappin}.} So the name of the fruit first meant the tool, then the
cluster, and only last the single berry.} The rest of the genus, several
dozen species, grows wild across the temperate Northern Hemisphere, and a
good number of them are American. For most of history this mattered to no
one in Europe.

Then, in the 1860s, vines in the south of France began to
die.\note{The cause was an insect smaller than a pinhead, now called
\emph{Daktulosphaira vitifoliae} but known to everyone by its older name,
phylloxera.\subnote{From Greek \emph{phyllon}, leaf, and \emph{xeros}, dry.
The name describes the damage the insect does to leaves, although the fatal
damage is to the roots.} It had come from North America, probably on vines
brought over by collectors. American vines had lived with it for a long time
and tolerated it. European vines had not.} Over the next decades the blight
moved through France and into the rest of Europe. The cure, found after much
argument, was to graft European vines onto American roots. Most of the
world's wine grapes still grow that way, a European plant standing on
American feet.\note{The exceptions are places the insect never reached or
cannot live, such as vineyards on pure sand.\subnote{Sand seems to hinder
the insect's movement through the soil. Why, exactly, is a question I have
not yet looked into.} Chile is the example usually given.}

The Americans, meanwhile, had grapes of their own. In the 1840s Ephraim Wall
Bull, in Concord, Massachusetts,\aside[concord]{The same Concord as Emerson
and Thoreau.\note{Thoreau was born there, in 1817.}} began planting seeds of a wild native species, \emph{Vitis
labrusca}, and selecting among the seedlings. The Concord grape came out of
that work.\note{Its flavour, which Europeans tend to call foxy, is the
flavour most Americans know as grape: the flavour of grape jelly and grape
soda. In 1869 Thomas Welch, a dentist in Vineland, New Jersey, heated
Concord juice to stop it from fermenting, and so made an unfermented wine
for church communion. For the town, see \xref{concord}.\subnote{Welch was a supporter of temperance. He was
applying a method that Louis Pasteur had worked out only a few years
earlier, in the mid-1860s.}}

In the Tokaj hills of north-eastern Hungary,\aside[tokay]{Older English
spells it Tokay.} the growers wait for a mould. Late in a good autumn, a fungus pierces
the skins of the ripe berries and the water inside evaporates. What remains
shrivels into something between a raisin and a jewel, and it makes one of
the great sweet wines.\note{The fungus is \emph{Botrytis cinerea}, noble
rot. The shrivelled berries are called \emph{aszú}.\subnote{The word is
related to the Hungarian verb for drying up. For the spelling of the place,
see \xref{tokay}.} In the wrong weather the same
fungus is simply grey mould, the grey fur on a forgotten punnet of
strawberries.\subnote[strawberry-returns]{Here the strawberry returns (see
\xref{strawberry}), as it will throughout this book, and with it the
question of what counts as a berry, \xref{not-berries}.}}
